\section{Introduction}
\label{sec:introduction}

Climate change has intensified the frequency and severity of extreme
weather events worldwide, with multiyear droughts becoming increasingly
common and vegetation resilience declining across multiple biomes
\citep{chen2025global,smith2023global}. Urban ecosystems are
particularly vulnerable because concentrated impervious surfaces,
fragmented green spaces, and anthropogenic heat amplification magnify
the impacts of heatwaves, droughts, and extreme precipitation
\citep{ipcc2023cities}. In response, national and subnational
governments have launched adaptation programmes designed to bolster
the capacity of urban ecosystems to withstand and recover from climate
disturbances. Evaluating whether such policies achieve their intended
ecological outcomes, however, has proven methodologically challenging:
randomised assignment is infeasible, treatment is endogenous to prior
environmental conditions, and conventional difference-in-differences
(DID) estimators impose functional-form assumptions that may mask
heterogeneous or nonlinear confounding relationships
\citep{rambachan2023more,chaisemartin2023twoway}.

China's climate-adaptive city pilot programme, launched in 2017 under
document No.~343 of the National Development and Reform Commission,
designated 28 prefecture-level cities as pilot sites for comprehensive
climate adaptation planning. The policy mandates integration of
climate-resilient infrastructure, ecological corridor construction,
and early-warning systems for extreme weather into municipal development
plans. Unlike the earlier sponge city programme (2015--2016), which
focused narrowly on stormwater management \citep{han2023chinas}, the
climate-adaptive city pilot encompasses a broader portfolio of
interventions targeting ecological resilience under multiple types of
climate stress. This broader scope, combined with the programme's
simultaneous city-level implementation, provides a quasi-experimental
setting well suited to causal evaluation.

Despite the policy's ambition, three gaps in the existing literature
limit our understanding of its effectiveness. First, most evaluations
of Chinese urban environmental pilots focus on economic outcomes such
as green total factor productivity or ESG performance
\citep{yuan2024double,qian2025impact}, while direct measurement of
ecological resilience---the policy's stated objective---remains
understudied. Second, studies that do examine ecological outcomes
typically rely on composite indices constructed from cross-sectional
data or simple trend analysis, without accounting for the role of
extreme weather events as the very stressors the policy is designed to
buffer against \citep{gong2025beyond,hasan2025analyzing}. Third, the
dominant empirical approach---two-way fixed effects (TWFE) DID---has
been shown to produce biased estimates under heterogeneous treatment
effects and staggered adoption \citep{borusyak2024revisiting,
wooldridge2025twoway}, motivating the adoption of more flexible
estimators that relax parametric assumptions about the
confounding structure.

This study addresses these gaps by making four contributions. First,
we construct a Climate-Stress Ecological Elasticity (CSEE) index that
directly measures ecosystem resistance to and recovery from extreme
weather events using satellite-derived NDVI, thereby aligning the
outcome variable with the policy's resilience objective. Second, we
employ a Double Machine Learning (DML) framework
\citep{chernozhukov2022automatic} that combines DID identification
with cross-fitted machine-learning estimation of nuisance functions,
relaxing the functional-form constraints of parametric DID while
preserving the causal interpretation of the treatment effect
\citep{ahrens2024ddml}. Third, we integrate multiple satellite data
sources---MODIS NDVI, ERA5 reanalysis, and MODIS land surface
temperature---with municipal yearbook data to build a rich panel of
338 cities over 2005--2023, achieving 77\% coverage with real
socioeconomic controls. Fourth, we conduct a comprehensive battery of
robustness tests, including 200-iteration placebo tests, propensity-score
matching, subsample restrictions, and an alternative-outcome analysis
using land surface temperature, alongside a seven-dimension
heterogeneity analysis and a resistance--recovery mechanism
decomposition.

The remainder of this paper is organised as follows.
Section~\ref{sec:literature} reviews the relevant literature on
ecological resilience measurement, climate adaptation policy
evaluation, and machine-learning causal inference.
Section~\ref{sec:methodology} introduces the DML--DID framework and
the CSEE construction. Section~\ref{sec:data} describes the data
sources and variable construction. Section~\ref{sec:results} presents
the main estimation results, parallel-trends tests, heterogeneity
analysis, and robustness checks. Section~\ref{sec:discussion} discusses
the findings, policy implications, and limitations.
Section~\ref{sec:conclusion} concludes.
